% ============================================================
%  BAB 1 – PENDAHULUAN
%  Compile standalone:  pdflatex bab-1.tex
%  Compile as part of thesis: pdflatex main.tex
% ============================================================
\documentclass[main]{subfiles}
\usepackage{hyphenat}
\usepackage{iftex} % Detects the compiler

% ── Encoding, Language & Fonts ───────────────────────────────
\ifXeTeX
    % Use this if compiling with XeLaTeX / LuaLaTeX
    \usepackage{polyglossia}
    \setmainlanguage{bahasai}   
    
    \usepackage{fontspec}
    \setmainfont{TeX Gyre Termes}   % Times New Roman equivalent
    \setsansfont{TeX Gyre Heros}    % Arial equivalent
    \setmonofont[Scale=0.88]{TeX Gyre Cursor}
\else
    % Use this if compiling with standard pdfLaTeX (Overleaf Default)
    \usepackage[utf8]{inputenc}
    \usepackage[T1]{fontenc}
    \usepackage[indonesian]{babel}
    
    \usepackage{mathptmx} % Times New Roman equivalent
    \usepackage{helvet}   % Arial equivalent
    \usepackage{courier}
\fi

% ── Page geometry (A4, UI thesis margins) ────────────────────
\usepackage[
  a4paper,
  top    = 2.54cm,
  bottom = 2.54cm,
  left   = 3.17cm,
  right  = 2.54cm
]{geometry}

% ── Line spacing & paragraph ─────────────────────────────────
\usepackage{setspace}
\onehalfspacing
\setlength{\parindent}{1.27cm}
\setlength{\parskip}{0pt}

% ── Better justification & hyphenation ───────────────────────
\usepackage{microtype}
\sloppy

% ── Headings (chapter/section styled like UI thesis) ─────────
\usepackage{titlesec}

% Chapter: centered, bold, uppercase, no "Chapter N" prefix
\titleformat{\chapter}[block]
  {\centering\sffamily\bfseries\large}
  {}
  {0pt}
  {\MakeUppercase}
\titlespacing*{\chapter}{0pt}{0pt}{1.5em}

% Section: left-aligned, bold, numbered
\titleformat{\section}
  {\sffamily\bfseries\normalsize}
  {\thesection}
  {1em}{}
\titlespacing*{\section}{0pt}{1.2em}{0.6em}

% Subsection
\titleformat{\subsection}
  {\sffamily\bfseries\normalsize}
  {\thesubsection}
  {1em}{}
\titlespacing*{\subsection}{0pt}{1em}{0.5em}

% ── List formatting ──────────────────────────────────────────
\usepackage{enumitem}
\setlist[enumerate, 1]{
  label        = \arabic*.,
  leftmargin   = 1.27cm,
  itemsep      = 4pt,
  parsep       = 0pt,
  topsep       = 4pt
}
\setlist[enumerate, 2]{
  label        = \alph*.,
  leftmargin   = 1.27cm,
  itemsep      = 3pt,
  parsep       = 0pt,
  topsep       = 2pt
}

% ── Hyperlinks (optional – comment out if submitting PDF/A) ──
\usepackage[hidelinks]{hyperref}

% Bibliography setup is inherited from main.tex (biblatex + references.bib)
% when this file is compiled standalone via the subfiles class.

% ============================================================
\begin{document}

\setcounter{chapter}{0}   % BAB 1

% ── Chapter heading ──────────────────────────────────────────
\chapter{Pendahuluan}

% ============================================================
\section{Latar Belakang}
% ============================================================

Kurikulum Merdeka, yang ditetapkan melalui Permendikbudristek Nomor 12 Tahun
2024, menempatkan \textit{meaningful learning} dan penguasaan kompetensi lintas
disiplin sebagai prioritas utama pendidikan nasional. Namun, pada jenjang SMA
khususnya Fase~F, rumpun Ilmu Pengetahuan Alam (IPA) tetap terstruktur menjadi
tiga mata pelajaran terpisah: Fisika, Kimia, dan Biologi.
Pemisahan ini berisiko menimbulkan fenomena \textit{siloed learning}
\parencite{harvey_silos_2025}, kondisi di mana batas-batas disipliner yang kaku
menghalangi peserta didik untuk mengonstruksi koneksi konseptual
antarmatapelajaran. Akibatnya, siswa kesulitan memahami keterkaitan yang secara
ontologis menghubungkan ketiga disiplin ilmu tersebut, sementara guru menghadapi
keterbatasan referensi kurikulum lintas-disiplin yang tersedia secara terstruktur.

Sebagai solusi atas fragmentasi struktur kurikulum ini, pendekatan berbasis
\textit{Knowledge Graph} (KG) menawarkan paradigma representasi yang mampu
memodelkan interkoneksi materi secara eksplisit. \textcite{ji_survey_2022} mendefinisikan
KG sebagai struktur data grafis yang merepresentasikan entitas dunia nyata dan
relasi semantiknya dalam format \textit{triple} (\textit{head, relation, tail}),
yang tidak hanya menyimpan fakta tetapi juga mendukung penalaran formal atas
data tersebut. \textcite{kim_study_2023} mengajukan pendekatan KG untuk
memvisualisasikan \textit{flow of knowledge} dan \textit{thinking style} yang
esensial bagi siswa guna membangun inferensi logis antar-konsep. Urgensi
penerapan ini diperkuat oleh \textcite{abu-salih_systematic_2024}, yang melalui tinjauan
sistematisnya menyimpulkan bahwa KG berperan vital dalam desain kurikulum dan
pemetaan konsep, khususnya untuk memvisualisasikan struktur pengetahuan yang
kompleks dan merancang jalur pembelajaran yang adaptif.

Meskipun menjanjikan, KG memiliki keterbatasan fundamental, yaitu sifatnya yang
secara inheren tidak lengkap. Banyak hubungan antar-konsep yang valid tidak
tertangkap ketika graf dibangun, baik melalui anotasi manual maupun ekstraksi
otomatis \parencite{ji_survey_2022}. Dalam konteks kurikulum sains, relasi lintas-disiplin
yang paling bernilai pedagogis justru seringkali bersifat implisit---tidak tersurat
dalam teks satu buku tunggal. Untuk mengatasi hal ini,
\textit{Knowledge Graph Completion} (KGC) diterapkan untuk memprediksi
\textit{triple} yang hilang. Akan tetapi, metode KGC konvensional berbasis
\textit{embedding} struktural terbukti tidak efektif dalam skenario kurikulum
karena distribusi konsep yang tidak merata. \textcite{wei_kicgpt_2023} menegaskan bahwa
metode berbasis \textit{triple} mengalami penurunan performa drastis saat
berhadapan dengan \textit{long-tail entities}---entitas yang memiliki sedikit
koneksi dalam graf sehingga pola struktural yang tersedia tidak cukup untuk
menghasilkan representasi vektor yang akurat.

Kemunculan \textit{Large Language Models} (LLM) membuka solusi yang menjanjikan
atas keterbatasan tersebut. \textcite{pan_unifying_2024} mengklasifikasikan pendekatan
LLM-augmented KGs melalui dua peran utama: sebagai \textit{encoder} yang
menghasilkan representasi vektor lebih kaya makna dibanding \textit{embedding}
struktural semata, serta sebagai \textit{generator} yang langsung memprediksi
\textit{triple} baru berdasarkan konteks. \textcite{yao_exploring_2025} mendemonstrasikan
bahwa dengan memperlakukan \textit{triple} pengetahuan sebagai sekuens teks,
LLM mampu memprediksi relasi semantik dengan akurasi yang lebih baik dari metode
konvensional. Secara khusus untuk mengatasi kelangkaan informasi pada entitas
\textit{long-tail}, \textcite{wei_kicgpt_2023} membuktikan bahwa strategi
\textit{In-Context Learning} (ICL)---yakni \textit{prompt few-shot} pada tahap KGC
yang menyertakan beberapa contoh \textit{triple} valid sebagai panduan di dalam
\textit{prompt} klasifikasi---memungkinkan LLM menginferensikan
relasi yang tidak mampu ditangkap oleh metode berbasis \textit{embedding}.

Pemanfaatan LLM dalam KGC tetap memerlukan validasi pakar
(\textit{human-in-the-loop}) untuk meminimalisir risiko \textit{AI hallucination}.
Meski demikian, evaluasi pakar secara manual kerap menimbulkan ketidaksepakatan
antar-evaluator yang subjektif dan sulit diterapkan pada skala besar
\parencite{gu_survey_nodate}. Penelitian ini mengatasi masalah tersebut dengan
mengadopsi pendekatan \textit{LLM-as-a-judge}, di mana LLM diposisikan sebagai
penengah yang berlandaskan konteks buku teks sumber untuk menyelesaikan konflik
penilaian antarpakar, sehingga validitas keilmuan KG yang dihasilkan dapat
dipertanggungjawabkan \parencite{gu_survey_nodate}.

Meskipun pendekatan \textit{LLM-Assisted Knowledge Graph Completion} telah
terbukti efektif pada domain umum, belum ada studi yang secara spesifik
mengadaptasinya untuk memetakan interkoneksi materi sains (Fisika, Kimia,
Biologi) dalam konteks Kurikulum Merdeka berbahasa Indonesia. Penelitian ini
mengisi kesenjangan tersebut dengan mengembangkan \textit{pipeline}
LLM-Assisted KGC untuk membangun dan melengkapi graf pengetahuan sains
lintas-disiplin berbasis buku teks resmi Kemendikbudristek.

% ============================================================
\section{Rumusan Masalah}
% ============================================================

Berdasarkan latar belakang di atas, penelitian ini merumuskan empat pertanyaan
penelitian berikut.

\begin{enumerate}
  \item \textbf{Masalah Fragmentasi Kurikulum:} Bagaimana memetakan
        interkoneksi semantik antar-materi sains (Fisika, Kimia, Biologi) pada
        Fase~F Kelas~XII Kurikulum Merdeka untuk mengatasi fenomena
        \textit{siloed learning} akibat struktur mata pelajaran yang terpisah?

  \item \textbf{Keterbatasan Metode Konvensional:} Bagaimana mengatasi kendala
        ketidaklengkapan dan penurunan performa metode KGC berbasis
        \textit{embedding} struktural saat berhadapan dengan entitas
        \textit{long-tail} dalam materi kurikulum yang distribusinya tidak merata?

  \item \textbf{Implementasi LLM-Assisted KGC:} Bagaimana efektivitas penerapan
        LLM-Assisted KGC dalam memprediksi relasi implisit lintas-disiplin
        antar-materi sains pada Kurikulum Merdeka?

  \item \textbf{Validasi Hybrid dan Evaluasi Metrik:} Bagaimana mengevaluasi
        keandalan dan kebenaran relasi sains yang dihasilkan LLM melalui
        \textit{pipeline} evaluasi hybrid (\textit{LLM-as-a-judge}) dan mengukur
        kualitas KG yang dihasilkan secara kuantitatif?
\end{enumerate}

% ============================================================
\section{Tujuan Penelitian}
% ============================================================

\subsection{Tujuan Umum}

Penelitian ini bertujuan menghasilkan purwarupa \textit{Knowledge Graph}
kurikulum sains Fase~F Kelas~XII yang memetakan interkoneksi lintas-disiplin
antara Fisika, Kimia, dan Biologi secara eksplisit dan tervalidasi. Graf yang
dihasilkan diharapkan menjadi fondasi data bagi pengembang kurikulum dan peneliti
teknologi pendidikan dalam merancang representasi pengetahuan yang mengakomodasi
keterkaitan antar-konsep lintas mata pelajaran.

\subsection{Tujuan Khusus}

Secara spesifik, penelitian ini bertujuan untuk:

\begin{enumerate}
  \item Membangun \textit{knowledge graph} awal dengan mengekstraksi entitas
        materi dan relasi eksplisit dari buku teks siswa Fisika, Kimia, dan
        Biologi Kelas~XII (e-book resmi Kemendikbudristek) menggunakan
        \textit{pipeline} LLM berbasis skema ontologi OWL, sebagai fondasi
        struktur pengetahuan kurikulum yang terstandarisasi.

  \item Mengimplementasikan mekanisme LLM-Assisted \textit{Knowledge Graph
        Completion} untuk memprediksi dan
        melengkapi \textit{missing links} lintas-disiplin---relasi implisit
        antar-entitas sains yang tidak tertangkap oleh ekstraksi berbasis aturan.

  \item Mengevaluasi kualitas \textit{knowledge graph} yang dihasilkan melalui
        dua pendekatan yang saling melengkapi:
        \begin{enumerate}
          \item \textbf{Evaluasi Struktural:} mengukur perubahan topologi graf
                sebelum dan sesudah proses KGC menggunakan metrik
                \textit{Average Degree Centrality} (ADC),
                \textit{Graph Density}, \textit{Modularity},
                \textit{Typed Relation Ratio}, dan
                \textit{Description Completeness}.

          \item \textbf{Validasi Pakar:} memverifikasi akurasi ontologis dan
                relevansi pedagogis relasi yang diprediksi LLM melalui enam
                \textit{reviewer} domain (dua per mata pelajaran), dengan metrik
                \textit{Fuzzy Precision and Recall},
                pengujian kesepakatan Gwet's AC1, dan penengahan ketidaksepakatan
                menggunakan \textit{pipeline LLM-as-a-judge}.
        \end{enumerate}
\end{enumerate}

% ============================================================
\section{Manfaat Penelitian}
% ============================================================

\subsection{Bagi Peneliti dan Pengembang Teknologi Pendidikan}

Penelitian ini menghasilkan dua kontribusi yang dapat dilanjutkan oleh peneliti
berikutnya. Pertama, \textit{pipeline} LLM-KGC yang didokumentasikan dalam
skripsi ini---mencakup skema ontologi OWL beserta strategi
\textit{prompt engineering}-nya---dapat diadaptasi untuk domain kurikulum lain
atau jenjang pendidikan yang berbeda. Kedua, KG kurikulum sains Fase~F Kelas XII berbahasa Indonesia
yang dihasilkan merupakan sumber daya terbuka pertama dari jenisnya, dan dapat
menjadi \textit{ground truth} awal untuk penelitian pembelajaran mesin di bidang
pendidikan nasional.

\subsection{Bagi Pendidik dan Pengembang Kurikulum}

Visualisasi KG yang dihasilkan menyediakan representasi eksplisit dari jembatan
konseptual lintas-disiplin yang selama ini bersifat implisit---misalnya
ketergantungan konsep Potensial Elektrode di Kimia terhadap Potensial Listrik di
Fisika, atau peran Reaksi Oksidasi sebagai prasyarat pemahaman Fosforilasi
Oksidatif di Biologi. Representasi ini dapat digunakan sebagai referensi ketika
merancang unit pembelajaran terpadu atau mengidentifikasi urutan konsep prasyarat
lintas mata pelajaran, tanpa harus membangun pemetaan tersebut secara manual.

% ============================================================
\section{Ruang Lingkup Penelitian}
% ============================================================

Berikut adalah batasan-batasan yang ditetapkan untuk penelitian ini.

\begin{enumerate}
  \item \textbf{Objek Penelitian:} Korpus data terbatas pada buku teks siswa
        (e-book PDF resmi Kemendikbudristek) mata pelajaran Fisika, Kimia, dan
        Biologi Fase~F Kelas~XII, diperoleh dari repositori
        \texttt{buku.kemdikbud.go.id}. Buku teks Kelas~XI dan dokumen
        Capaian Pembelajaran (CP) tidak termasuk dalam korpus.

  \item \textbf{Model dan Strategi LLM:} Model Gemini~2.5 Flash Lite
        (Google GenAI SDK) digunakan untuk ekstraksi entitas maupun
        penyempurnaan graf (\textit{completion}) relasi lintas-disiplin, sedangkan
        Gemini~3.1 Pro Preview digunakan khusus sebagai \textit{LLM-as-a-judge}
        dalam \textit{pipeline} validasi.
        \textit{Output} seluruh model ditetapkan dalam Bahasa Indonesia.
        Penelitian ini tidak mencakup \textit{pre-training} atau
        \textit{fine-tuning} model baru; LLM dimanfaatkan melalui
        \textit{prompt engineering} berbasis \textit{few-shot} dan
        \textit{schema-constrained}, baik untuk ekstraksi entitas maupun untuk
        klasifikasi relasi lintas-disiplin pada tahap \textit{completion}.

  \item \textbf{Infrastruktur Graf:} Graf diinstansiasi pada basis data Neo4j
        (Aura Free tier) menggunakan Cypher. Pencarian kemiripan semantik
        memanfaatkan ANN \textit{vector index} HNSW dengan model
        \textit{embedding} \texttt{gemini-embedding-001}. Metrik topologi (ADC,
        \textit{Graph Density}, \textit{Modularity}) dihitung melalui kueri
        Cypher.

  \item \textbf{Cakupan Validasi:} Validasi pakar dilakukan oleh enam
        \textit{reviewer} (dua per mata pelajaran) dengan latar belakang ilmu
        murni atau pendidikan sains yang linier. Setiap \textit{triple} dinilai
        ke dalam empat kategori penilaian melalui aplikasi web KG Review App. Kesepakatan
        diukur menggunakan Gwet's AC1. Ketidaksepakatan pada validasi
        intra-mapel diselesaikan oleh \textit{LLM-as-a-judge} berbasis
        Gemini~3.1 Pro Preview; ketidaksepakatan pada relasi lintas-disiplin
        diselesaikan melalui diskusi panel seluruh pakar.

  \item \textbf{Luaran dan Batasan Sistem:} Luaran penelitian adalah purwarupa
        struktur data \textit{Knowledge Graph} pada Neo4j beserta visualisasinya.
        Penelitian ini tidak mencakup pengembangan sistem rekomendasi pembelajaran
        adaptif, antarmuka pengguna berbasis web (\textit{full-stack
        application}), atau integrasi ke dalam platform LMS.
\end{enumerate}

\ifSubfilesClassLoaded{\printbibliography[title={Daftar Pustaka}]}{}

\end{document}